\section{从太原到洛阳：建国不是一日的受禅}
\label{sec:early-tang-foundation}

617年李渊自太原起兵，入关后拥立代王杨侑并受封唐王；次年受禅建唐。\eventanchor{ST-617-LI-YUAN-TAIYUAN}{太原起兵（617）}太原军镇、关陇网络和与突厥的交涉，使他能够取得长安，却不等于已经取得一个可直接调度的帝国。唐初史书特别突出李世民的作用和突厥关系，正是后来的继承与军功记忆塑造了材料；可以把握的是关中出现了一个能接收隋朝官署与财政名义的中心。%
\footnote{太原起兵与唐建国见 ST-617-LI-YUAN-TAIYUAN、ST-618-TANG-FOUNDING，依据《旧唐书·高祖本纪》《新唐书·高祖本纪》及《资治通鉴》隋、唐纪。结盟条款和个人作用的细节受唐代叙事重塑。}

\begin{event}{太原起兵与建唐：617--618年，关中}
起兵、入关和受禅构成可相互校读的年序；它们并不证明新朝已经在各地取得同样的资源、服从或行政能力。后文循官署、道路与军政网络追问的，正是这项转换如何被继续完成。
\end{event}

王世充在洛阳建郑，窦建德的夏军又能自河北南援；621年虎牢之战及王世充降唐，使唐取得洛阳与中原仓储，却没有终结地方接管。次年刘黑闼在河北再起，说明军队、旧部和地方征发的关系不能由一纸降书消除。唐朝的统一战争因此不是一路“平定”，而是围绕洛阳、河北道路与州县人事反复建立可用控制。%
\footnote{郑建国、虎牢与河北复叛见 ST-619-ZHENG-FOUNDING、ST-621-HULAO、ST-622-LIU-HEITA，依据《旧唐书》王世充、窦建德、刘黑闼诸传、《新唐书》相关列传及《资治通鉴》。兵数、突阵和地方支持度主要来自胜者的军功叙事。}

626年玄武门政变将继承争夺置于禁军和宫城的暴力中；两月后颉利可汗至渭水北岸，唐廷以会盟换得退兵。二者相接，正说明新皇帝不能把内部继承与边防分开处理。628年梁师都政权覆灭，才使关中北缘较为稳固；唐的早期国家能力是在军队、任官、粮储与边地外交都尚未定型时逐步取得的，而不是贞观一开始便自然完备。%
\footnote{玄武门、渭水与梁师都覆亡见 ST-626-XUANWU、ST-626-WEISHUI、ST-628-LIANG-SHIDU，依据《旧唐书·太宗本纪上》《旧唐书·突厥传上》《旧唐书·梁师都传》及《资治通鉴》。政变预谋和渭水谈判细节均有后出的修辞层次。}

\subsection{交叉阅读窗口二：受禅、律文与滕王阁的远望}

617--628年的建国叙事，宜以《旧唐书》《新唐书》高祖、太宗纪和《资治通鉴》的年序相互牵制：前者留下唐朝开国者及北宋重修者的不同取景，后者将兵变、受禅与边事压成因果链。永徽四年（653）颁定的《唐律疏议》与后出《唐会要》所辑制度材料，不能倒过来证明621年每一州县已经依同一格式裁判；它们却使我们能问，唐为何要把隋以来的官署、户婚和刑名转化为可复用的行政语言。初唐墓志所见籍贯、官衔与婚姻可作局部校验，但其作者、受葬者皆有选择，不能补成统一战争中普通家庭的档案。%
\footnote{建国与统一战争可互校《旧唐书·高祖本纪》《旧唐书·太宗本纪上》《新唐书》相关本纪、列传与《资治通鉴》唐纪；律文见《唐律疏议》，制度条目可参《唐会要》。永徽律的年代晚于本节多数事件，故这里只把它当作制度延续的提问，不当作617--628年的直接证词。}

王勃《滕王阁序》有“关山难越，谁悲失路之人；萍水相逢，尽是他乡之客”之句。此篇为骈文宴集之作，传统系于675年前后，作者以早慧而失意的宾客身份铺陈江山、仕途与相逢；它的句法把交通节点的壮阔转成个人去留的压力。它不能证明唐初某条道路的人口流向，时间也晚于虎牢和渭水数十年；但它提醒读者，统一不是把关山抹平，而是让道路、州县名分和仕进机会重新决定谁能够抵达中心。对钱穆所问中央制度如何获得地方支点而言，这是一种可检验的空间问题；对吕思勉的社会史问题而言，则须回到律文、墓志和文书，而不能让一篇华丽骈文代替迁徙史。
